For each prime form $H_p$, the set of residues reached by $2,2^2,2^3,\ldots$ modulo $p$, stopping when a residue repeats. Then
\[N_p=\#\{1\le n<p:\ (-n^2-1)\bmod p\in H_p\}.\]
This finite test is exhaustive: after a residue repeats, multiplying by 2 continues the same cycle. For odd $p$, one can shorten the count by checking $n=1,\ldots,\dfrac{p-1}{2}$ and doubling, because $n$ and $p-n$ have the same square. For $p=2$, direct checking gives $N_2=1$.

The full table is
\[\begin{array}{c|r@{\qquad}c|r@{\qquad}c|r}
p&N_p&p&N_p&p&N_p\\\hline
2&1&29&26&61&58\\
3&2&31&4&67&66\\
5&2&37&34&71&36\\
7&4&41&18&73&10\\
11&10&43&16&79&40\\
13&10&47&24&83&82\\
17&6&53&50&89&14\\
19&18&59&58&97&46\\
23&12&&&&
\end{array}\]
For example, $H_7=\{1,2,4\}$, and the admissible $n$ are $2,3,4,5$. The qualifying primes are $7,23,31,43,47,71,79$, whose sum is $\boxed{301}$.
